\documentclass[12pt,a4paper]{article}

\usepackage[utf8]{inputenc}
\usepackage[T1]{fontenc}
\usepackage[spanish]{babel}
\usepackage{geometry}
\usepackage{setspace}
\usepackage{csquotes}
\usepackage{hyperref}

\geometry{margin=2.5cm}
\onehalfspacing

\title{Del crédito, del visado y del crédito del visado}
\author{Vicente Domínguez Arca}
\date{}

\begin{document}

\maketitle

\section*{Introducción}

Resulta difícil objetar que el pasaporte de una persona constituye hoy uno de los múltiples elementos que intervienen en la producción y reproducción de la desigualdad global. No porque el pasaporte, por sí solo, agote la explicación de dicha desigualdad, sino porque distribuye de forma diferencial una facultad decisiva: la posibilidad de moverse, residir, trabajar, solicitar protección, acceder a mercados laborales o simplemente atravesar fronteras sin ser sometido a una sospecha administrativa previa.

El pasaporte suele presentarse como un documento neutral de identificación. Sin embargo, desde una perspectiva de teoría crítica, puede ser comprendido como un dispositivo institucional que condensa pertenencia política, reconocimiento jurídico, confianza estatal y valor geopolítico. En él se cruzan el derecho, la soberanía, la administración, la economía y la moral pública de las fronteras.

El propósito de este texto no es elaborar un tratado moral sobre la migración, ni reducir el fenómeno migratorio a una consecuencia lineal del capitalismo. La movilidad humana responde a causas múltiples: guerras, persecuciones, catástrofes, expectativas de vida, vínculos familiares, desigualdades económicas, transformaciones laborales y deseos legítimos de mejora personal. Sin embargo, una crítica seria no puede ignorar que, en el orden contemporáneo, el crédito económico previo puede operar como condición suficiente para obtener el ``crédito'' de numerosos Estados: visados, permisos de residencia e incluso, en ciertos casos, pasaportes.

\section*{Genealogía: pasaporte, paso y fiscalidad}

La historia del pasaporte no posee un origen único. Existen múltiples antecedentes antiguos y medievales: salvoconductos, cartas de tránsito, documentos de protección, permisos de circulación y certificaciones emitidas por autoridades políticas o religiosas. Entre esos antecedentes aparece con frecuencia la \textit{bara'a}, utilizada en determinados contextos islámicos medievales como certificación de que una persona había cumplido sus obligaciones fiscales y no mantenía deudas con el erario.

No se trata, por tanto, de afirmar que la \textit{bara'a} sea el origen exclusivo del pasaporte moderno. Más bien debe entenderse como uno de los ejemplos antiguos que dan inicio, junto con otros, a una misma lógica institucional: para desplazarse legítimamente, el sujeto debía ser reconocido por una autoridad y, en muchos casos, debía acreditar que no dejaba tras de sí una deuda fiscal, jurídica o política.

Este punto es decisivo. La movilidad no aparece históricamente como una pura libertad natural, sino como una libertad condicionada por formas de autorización. El desplazamiento requiere paso; el paso requiere permiso; el permiso requiere reconocimiento; y el reconocimiento suele apoyarse en algún tipo de crédito.

Incluso la etimología de ``pasaporte'' conserva esta ambigüedad material. El término francés \textit{passeport} remite a la acción de pasar y a un \textit{port} que puede entenderse, según la tradición etimológica considerada, como puerto, puerta, paso o lugar de tránsito. Lo relevante aquí no es resolver definitivamente la cuestión filológica, sino advertir que el término conserva la imagen de una barrera: algo que debe ser cruzado mediante autorización.

\section*{Tres sentidos del crédito}

La hipótesis central de este texto exige distinguir tres sentidos de crédito.

En primer lugar, crédito como confianza. Procede del latín \textit{credere}: creer, confiar. Todo pasaporte es una tecnología de credibilidad. El Estado afirma: esta persona es quien dice ser, pertenece a esta comunidad política y debe ser reconocida como tal.

En segundo lugar, crédito como regularidad fiscal o administrativa. En antecedentes como la \textit{bara'a}, el documento no sólo identificaba, sino que acreditaba una situación de cumplimiento: el sujeto no debía al erario, no escapaba de una obligación pendiente, no abandonaba una deuda pública.

En tercer lugar, crédito como capacidad financiera. Este es el sentido contemporáneo más explícito: el sujeto con recursos, patrimonio, liquidez, avales o capacidad de inversión puede presentarse ante el Estado receptor como alguien económicamente deseable. Su capital previo se convierte en signo de confiabilidad futura.

La tesis crítica surge de la superposición de estos tres sentidos: el crédito financiero puede transformarse en crédito institucional. Quien posee recursos económicos no sólo puede comprar bienes o servicios; en determinados marcos jurídicos, puede comprar confianza, reducir sospechas, acelerar procedimientos y acceder a formas reforzadas de movilidad.

\section*{El pasaporte como capital institucional heredado}

El pasaporte moderno acredita formalmente la nacionalidad. Pero en la práctica representa también una posición dentro de la jerarquía internacional de los Estados. No todos los pasaportes valen lo mismo. Algunos abren fronteras; otros las cierran. Algunos producen confianza inmediata; otros activan sospecha, demora, entrevista, prueba documental y coste económico.

En este sentido, el pasaporte puede entenderse como una forma de capital institucional heredado. No se elige el lugar de nacimiento, pero el lugar de nacimiento determina inicialmente el campo de posibilidades jurídicas de una persona. La nacionalidad funciona así como una renta de posición: quien nace bajo un pasaporte fuerte hereda una capacidad de movimiento que no ha producido individualmente.

Aquí resulta pertinente el diálogo con Bourdieu. Si el capital cultural o social permite ocupar posiciones diferenciales dentro de un campo, el pasaporte opera como una modalidad de capital político-institucional. No es simplemente un documento; es una credencial de acceso a oportunidades desigualmente distribuidas.

\section*{El Estado y los medios legítimos de movimiento}

John Torpey ha mostrado que el Estado moderno no sólo monopoliza los medios legítimos de violencia, como afirmaba Weber, sino también los medios legítimos de movimiento. El pasaporte, el visado, el registro civil, la identificación personal y la frontera administrativa forman parte de esa apropiación estatal de la movilidad.

Desde esta perspectiva, el pasaporte no es un residuo burocrático menor, sino una pieza central de la soberanía moderna. El Estado identifica, clasifica, autoriza, excluye y documenta. Decide quién puede circular, bajo qué condiciones, durante cuánto tiempo y con qué derechos.

La aportación específica de este texto consiste en añadir que esa monopolización de la movilidad está atravesada por una lógica crediticia. El Estado no sólo pregunta quién eres, sino qué confianza mereces. Y esa confianza puede formularse jurídicamente, policialmente, laboralmente, culturalmente o económicamente.

\section*{El visado como evaluación crediticia}

El visado constituye una autorización adicional. Mientras el pasaporte acredita una pertenencia política de origen, el visado expresa una evaluación del Estado receptor. El solicitante debe demostrar que no representa una amenaza, que posee medios suficientes, que tiene motivos legítimos de entrada, que no permanecerá irregularmente o que aportará algún valor al país receptor.

En ese sentido, el visado funciona como una evaluación crediticia ampliada. No se limita al crédito bancario, aunque puede incluirlo. Evalúa solvencia económica, trayectoria profesional, formación, estabilidad familiar, vínculos de retorno, historial jurídico y expectativas de integración.

El migrante cualificado, con ingresos, contrato, ahorros, patrimonio o una profesión demandada, recibe más fácilmente el crédito institucional del Estado receptor. El migrante pobre, incluso cuando migra por razones comprensibles o legítimas, debe demostrar mucho más. No basta con querer pasar; debe convencer de que merece pasar.

\section*{Ciudadanía por inversión: el crédito que compra crédito}

Esta lógica alcanza su forma extrema en los programas de residencia o ciudadanía por inversión. En ellos, determinados Estados permiten obtener residencia, e incluso nacionalidad, a cambio de inversiones, donaciones, compra de deuda pública, adquisición inmobiliaria u otros compromisos económicos.

No conviene convertir este fenómeno en el centro absoluto del análisis, pues no representa la experiencia migratoria común. Sin embargo, una crítica seria no puede obviarlo. Estos programas muestran con claridad que el crédito previo puede transformarse directamente en crédito político. El sujeto no sólo adquiere una casa, una participación económica o una posición financiera; puede adquirir también movilidad, residencia, ciudadanía e incluso un nuevo pasaporte.

Naturalmente, estos programas suelen incluir mecanismos de diligencia debida: comprobación de antecedentes, origen lícito de fondos, reputación pública y cumplimiento fiscal. Pero la existencia de tales controles no elimina la cuestión crítica. Al contrario, la confirma: el Estado convierte la ciudadanía en un objeto condicionado por una evaluación económico-moral del solicitante.

En sus versiones más problemáticas, estos programas han permitido que individuos con historiales controvertidos o con riesgos jurídicos significativos accedan a nuevas formas de movilidad internacional. No es lo normal, ni debe presentarse como regla general. Pero su mera posibilidad revela una asimetría fundamental: allí donde el migrante pobre debe probar necesidad, arraigo, obediencia y utilidad, el migrante rico puede presentar inversión.

\section*{Migración, formación y rentabilidad}

La cuestión se vuelve más compleja cuando se observa la migración ordinaria. Una persona puede traer experiencia laboral, capacidades prácticas y voluntad de integración, pero si sus habilidades no son inmediatamente traducibles al mercado receptor, aparece un coste de adaptación.

Ese coste puede adoptar varias formas: aprendizaje lingüístico, homologación de títulos, formación técnica, acompañamiento administrativo, vivienda, salud, escolarización o inserción laboral. Si el migrante posee crédito económico, puede financiar parte de ese tránsito. Si no lo posee, dependerá de redes familiares, comunitarias, públicas o empresariales.

Aquí aparece una tensión central del capitalismo contemporáneo. Las empresas pueden invertir en formación, pero tienden a hacerlo cuando esperan retorno. El Estado de bienestar puede asumir esa inversión, pero entonces el coste se socializa mediante gasto público y fiscalidad. El individuo puede intentar financiarse, pero sin avales ni historial crediticio local difícilmente obtendrá crédito bancario.

La paradoja es evidente: quien más necesita crédito para integrarse es, a menudo, quien menos crédito recibe. La falta de crédito económico se convierte en falta de crédito institucional, y la falta de crédito institucional refuerza la exclusión económica.

\section*{La tasa geopolítica del pasaporte}

Puede proponerse aquí una analogía con la usura. En el crédito financiero, el interés funciona como precio del acceso al dinero. En la movilidad global, el visado funciona en ocasiones como precio del acceso al territorio. Pero ese precio no es igual para todos. Depende del pasaporte de origen, del capital económico, de la profesión, de la raza administrativamente percibida, de la lengua, de la religión, del país emisor y de la posición geopolítica del solicitante.

El sistema internacional aplica así una suerte de tasa geopolítica diferencial. Algunos sujetos cruzan casi sin coste; otros pagan con dinero, espera, sospecha, humillación documental, entrevistas, traducciones, certificados y riesgo de rechazo.

El pasaporte fuerte reduce el interés de entrada. El pasaporte débil lo aumenta. El crédito económico puede compensar parcialmente esa debilidad. La pobreza, por el contrario, multiplica la sospecha.

\section*{Dimensión normativa}

Desde una teoría de la justicia, el pasaporte plantea un problema evidente: distribuye oportunidades vitales a partir de una circunstancia no elegida. En términos del igualitarismo de la suerte, la nacionalidad de nacimiento pertenece al ámbito de la fortuna bruta, no al mérito individual. Nadie merece moralmente un pasaporte fuerte del mismo modo que nadie merece moralmente un pasaporte débil.

Desde el enfoque de las capacidades, la movilidad no es un lujo secundario, sino una condición que puede afectar al trabajo, a la seguridad, a la vida familiar, a la salud, a la educación y al desarrollo personal. Limitar la movilidad de forma desigual significa limitar de forma desigual la posibilidad efectiva de vivir una vida elegida.

No se sigue de aquí que toda frontera sea ilegítima ni que todo control migratorio sea moralmente inadmisible. La cuestión es más precisa: si los Estados regulan la movilidad, deben poder justificar por qué ciertas personas heredan libertad de movimiento mientras otras heredan sospecha.

\section*{Conclusión}

El pasaporte y el visado no son meros instrumentos administrativos. Son dispositivos de reconocimiento, clasificación y selección. Acreditan identidad, pero también distribuyen confianza. Permiten pasar, pero también producen barreras. Organizan la movilidad, pero también jerarquizan vidas.

La noción de ``crédito del visado'' permite nombrar esta operación. El Estado receptor concede o niega crédito institucional al sujeto que solicita entrar. Ese crédito puede fundarse en la nacionalidad, la profesión, la solvencia, la historia fiscal, la reputación, la ausencia de antecedentes o la capacidad de inversión.

En el extremo, el crédito financiero puede convertirse en pasaporte. En la normalidad cotidiana, la falta de crédito económico dificulta obtener crédito institucional. Entre ambos polos se despliega una economía política de la movilidad.

Por ello, una crítica contemporánea del crédito no debe limitarse al banco, al préstamo o al interés. Debe atender también a las formas en que el crédito organiza la confianza social, la movilidad humana y la pertenencia política. El pasaporte es, en este sentido, una de las expresiones más visibles de una desigualdad heredada que se presenta como simple documentación.

\section*{Referencias orientativas}

\begin{thebibliography}{9}

\bibitem{torpey}
Torpey, John.
\textit{The Invention of the Passport: Surveillance, Citizenship and the State}. 2000.

\bibitem{bourdieu}
Bourdieu, Pierre.
\textit{Poder, derecho y clases sociales}. 2000.

\bibitem{foucault}
Foucault, Michel.
\textit{Seguridad, territorio, población}. 1978.

\bibitem{carens}
Carens, Joseph H.
\textit{The Ethics of Immigration}. 2013.

\bibitem{nussbaum}
Nussbaum, Martha C.
\textit{Crear capacidades: propuesta para el desarrollo humano}. 2012.

\bibitem{graeber}
Graeber, David.
\textit{En deuda: una historia alternativa de la economía}. 2012.

\bibitem{shachar}
Shachar, Ayelet.
\textit{The Birthright Lottery: Citizenship and Global Inequality}. 2009.

\bibitem{cohen}
Cohen, G. A.
\textit{On the Currency of Egalitarian Justice, and Other Essays in Political Philosophy}. 2011.

\end{thebibliography}

\end{document}
```
